\chapter{Bevezetés}
Szakdolgozatom témájaként a WaccApp nevű, WhatsApp Business API-ra épülő kommunikációs és kérdőívkezelő rendszer frontendjének fejlesztését választottam. A projekt alapötlete egy gyakorlati problémából indult ki: hogyan lehet egy technikailag összetett üzenetküldő platformot olyan webes felületté alakítani, amelyet a felhasználó nem API-hívások és háttérfolyamatok sorozataként érzékel, hanem egy átlátható, lépésekre bontható munkafolyamatként. 

A WhatsApp Business API jó alapot biztosít üzleti célú üzenetküldéshez, sablonos kommunikációhoz, interaktív válaszokhoz és médiatartalmak továbbításához. Ezek a lehetőségek azonban önmagukban még nem jelentenek könnyen használható rendszert. A fejlesztés során ezért nem csak az volt a feladatom, hogy a backend által biztosított funkciók elérhetők legyenek a böngészőből, hanem az is, hogy ezek a műveletek a felhasználó számára követhető és lehetőleg hibamentes folyamatként jelenjenek meg. [R1], [R3] 

A fejlesztés elején még úgy tűnt, hogy az üzenetküldés lesz a rendszer legfontosabb és legösszetettebb része. A gyakorlati próba közben viszont hamar kiderült, hogy önmagában az üzenet elküldése kevés. Amint megjelentek a sablonok, a kérdőívek és a visszakeresési igények, már nem egy egyszerű küldőfelületben kellett gondolkodnom, hanem egy olyan frontendben, amely több egymásra épülő munkafolyamatot tud kezelni.

Ez a bővülés a frontend szerkezetére is hatással volt. Egy idő után már nem lett volna célszerű minden funkciót egyetlen oldalra zsúfolni, mert az nehezen áttekinthető felületet eredményezett volna. Emiatt a rendszer több különálló, de egymáshoz kapcsolódó nézetre tagolódott. Az \path{index.html} lett a központi indítófelület, ahonnan az egyedi üzenetküldés, a sablonküldés, a kérdőívindítás, a médiaküldés és az időzítés is elérhető. A beszélgetések megtekintésére külön nézet szolgál, például a \path{conv.html}, míg a korábbi adatok keresését és szűrését a \path{filter.html} támogatja. 

A kérdőívek és sablonos folyamatok kezelése szintén külön felületet igényelt. Ezeknél már nem egyszerű adatbevitelről van szó, hanem olyan kommunikációs logikáról, ahol egy válasz befolyásolhatja a következő lépést. A frontendnek ezért nemcsak mezőket és gombokat kellett megjelenítenie, hanem azt is kezelnie kellett, hogy a felhasználó milyen sorrendben ad meg adatokat, milyen hibák fordulhatnak elő, és hogyan kap visszajelzést a művelet eredményéről. 

A téma aktualitását az adja, hogy az ügyfélkapcsolati, marketing- és információgyűjtési folyamatokban egyre fontosabbak azok a rendszerek, amelyek gyorsan, automatizálhatóan, mégis személyesebb formában tudnak kommunikálni a felhasználókkal. Egy ilyen rendszerben viszont nem elég, ha a háttérben működő technológia képes üzeneteket továbbítani. A kezelőfelületnek is támogatnia kell a felhasználót abban, hogy tudja, milyen adatot kell megadnia, mikor indítható el egy művelet, miért nem fogadható el egy hiányos kérés, és hogyan ellenőrizhető később egy elküldött üzenet vagy kérdőíves válasz. [R1], [R3] 

A dolgozat középpontjában ezért a frontendoldali megvalósítás áll. A backend és a WhatsApp Business API működése csak olyan mértékben jelenik meg, amennyire az a kliensoldali működés megértéséhez szükséges. A fejlesztés során kiemelt szerepet kapott a több HTML-nézetből álló felépítés, a reszponzív megjelenés, a kliensoldali validáció, az aszinkron API-kommunikáció kezelése, valamint a felhasználói visszajelzések következetes alkalmazása. 

A fejlesztés során több olyan helyzet is előfordult, amely egyértelműen befolyásolta a rendszer felépítését. Egy egyszerű szöveges üzenet elküldésénél például elegendő a címzett és az üzenetszöveg ellenőrzése. Egy sablonos kérdőív vagy időzített küldés viszont már több adatot igényel: sablont, paramétereket, válaszlehetőségeket és időpontot is kezelni kell. A kérdőívszerkesztésnél külön kihívást jelentett, hogy ne csak kérdéseket lehessen rögzíteni, hanem azt is, hogy egy adott válasz után melyik következő kérdés vagy sablon következzen. A szűrési felületen pedig az egyszerű listázás kevésnek bizonyult, mert a felhasználónak több szempont alapján is vissza kellett tudnia keresni az adatokat, például név, telefonszám, üzenettípus, dátumintervallum vagy kérdőíves válasz alapján. 

A szakdolgozat célja annak bemutatása, hogy a WaccApp frontendje hogyan rendezi webes kezelőfelületbe a WhatsApp Business API-ra épülő kommunikációs funkciókat. A dolgozat nem egy általános chatalkalmazást mutat be, hanem egy saját fejlesztésű frontend felépítését és működését. Ennek része az üzenetküldő felület, a beszélgetésnézet, a sablonos és szöveges kérdőívkezelés, a média- és fájlküldés, az időzített küldések kezelése, valamint a szűrési és egyszerű válaszösszesítési funkciók bemutatása.

A dolgozat elsőként az üzleti üzenetküldő rendszerek és a kapcsolódó frontendminták hátterét tekinti át. Ezt követően bemutatja a WaccApp technológiai alapjait, követelményeit és a felhasználói felület kialakításához kapcsolódó döntéseket. A későbbi fejezetekben a rendszer architektúrája, fő moduljai, API-kapcsolatai, kliensoldali biztonsági szempontjai és tesztelési tapasztalatai kerülnek bemutatásra. A dolgozat végén a megvalósított eredmények, a rendszer korlátai és a továbbfejlesztési lehetőségek összegzése szerepel. 